\section{模型的改进与推广}

\subsection{模型的改进}

后续改进首先需要补充多年销量、价格和亩产量数据。获得时间序列后，可按作物估计边际分布，并采用滚动时间窗进行样本外回测，以替代当前由题面区间设定的模拟分布。问题三还需要采购订单或价格联动数据，用于估计交叉价格弹性并检验关系边方向，避免把预设相关结构误当成真实市场机制。

在管理层面，可采集地块坐标、机械切换时间和分月劳动力上限，把运输与集中作业成本加入目标函数。新增约束会改变当前可行域，必须重新检查正式面积表及跨年轮作边界。若测试情景数量增加，可比较情景削减与分解方法的计算代价，但不能省略逐地块硬约束和样本外复核。

\subsection{模型的推广}

本文的面积变量、土地约束、轮作约束和分段销售函数可作为其他多地块、多品种、多周期配置问题的通用建模框架。迁移到设施农业、订单农业或区域种植结构调整时，需要重新提供地块适宜性、产销参数、轮作集合和销售规则。本文没有在其他地区或数据集上进行实验，因此这种推广只表示模型结构可重新配置，不表示现有方案具有跨区域有效性。

样本外评价框架还可以用于比较保险、最低收购价和合同销售等政策情景。此类推广只能在重新确认概率假设、政策约束和风险指标定义后进行，不能直接把本文的收益、亏损比例或宏观稳定结论外推到新的区域和政策环境。
